\chapter{Einleitung}
\label{sec:Einleitung}

Die vorliegende Untersuchung befasst sich mit der numerischen Simulation einer rechteckigen Scheibe, die durch eine zentrale elliptische Aussparung geschwächt wird. Die Geometrie ist durch eine horizontale Länge von ${l = 100\,\mathrm{mm}}$ und eine vertikale Höhe von ${h = 40\,\mathrm{mm}}$ charakterisiert, wobei die zentrale Ellipse Halbachsen von ${a = 6\,\mathrm{mm}}$ und ${b = 3\,\mathrm{mm}}$ aufweist. Um die mechanische Antwort unter einachsiger Belastung zu analysieren, wird das Modell an der linken Stirnfläche in ${x}$-Richtung eingespannt und an der gegenüberliegenden Seite mit einer Zugspannung von ${\sigma_0 = 100\,\mathrm{MPa}}$ beaufschlagt. Weiter wird der Freiheitsgrad in ${y}$-Richtung dir Lochscheibe an der Unterseite eingeschränkt.

Methodisch liegt der Fokus der Problemlösung auf einer Diskretisierungsstrategie, um trotz einer vorgegebenen Knotenobergrenze eine hinreichende Ergebnisgüte zu erzielen. Hierzu wird die Geometrie nach initialer Analyse mit einem unsortierten Netz zunächst gezielt in vier Quadranten partitioniert, um die Anwendung eines strukturierten Netzes zu ermöglichen. Das Vorgehen bei der Netzgenerierung konzentriert sich primär auf die lokale Verfeinerung am Ort der erwarteten Spannungskonzentration. Durch den Einsatz eines einseitigen Gradienten \texttt{Single Bias} wird die Knotendichte am Lochrand signifikant erhöht, während das Netz zu den Außenrändern hin sukzessive vergröbert wird. Ein wesentlicher Teil der Untersuchung widmet sich dabei der gezielten Auswahl der Elementformulierung. Hierbei wird insbesondere analysiert, inwieweit voll integrierte quadratische Elemente \texttt{CPS8} im Vergleich zu linearen Ansätzen geeignet sind, die komplexen Spannungsgradienten an der Kerbe präzise abzubilden und so eine Annäherung an die theoretischen Referenzwerte zu gewährleisten.